\chapter{Observed diffusion geometry and a compatible-model minimum}
\label{ch:og-population}

An observed continuous path contains local information that a list of
correlations need not retain. In a diffusion, quadratic variation records the
noise strength along the observed coordinate. A conditional comparison of
past and future increments then separates a directed contribution from
variation of the reversible drift inside an unresolved observation fiber.
The result below bounds entropy production without fixing the dimension or
architecture of the hidden diffusion. For one complete scalar process law,
the bound determines and attains the minimum over the entire declared hidden
model class.

These are population statements for ideal noiseless continuous observations.
They do not certify that a finite dataset has a specified process law. Their
physical interpretation requires ordinary reversal of the complete hidden
state; in particular, an equilibrium inertial model is not covered merely
because its measured position is even under reversal. The proof uses standard
diffusion calculus, conditional expectation, and the stationary reverse
semigroup. Historical priority for the combination and its consequences is
not asserted.

\section{Model class, entropy convention, and observed geometry}

Let $M$ be a compact connected smooth Riemannian manifold without boundary,
of arbitrary finite dimension. Let $\pi$ be a strictly positive smooth
probability density relative to its volume form, $D$ a smooth symmetric
uniformly positive definite diffusion tensor, and $v$ a smooth vector field
such that $\operatorname{div}(\pi v)=0$. On smooth real functions define
\begin{equation}
 L=S+A,\qquad
 Sf=\pi^{-1}\operatorname{div}(\pi D\nabla f),\qquad
 Af=v\mathbin{\cdot}\nabla f.
 \label{eq:og-generator}
\end{equation}
Consider the stationary two-sided diffusion $X=(X_t)_{t\in\mathbb R}$ with
generator $L$ and invariant density $\pi$. These assumptions imply
$L^\dagger=S-A$ in $L^2(\pi)$ and smooth bounded generator derivatives of
every smooth detector. Its stationary reverse process has generator
$L^\dagger$.

The convention for the diffusion coefficient is local covariance $2D\,dt$.
Write
\begin{equation}
 e(x)=v(x)^{\mathsf T}D(x)^{-1}v(x),\qquad
 \sigma=\int_M e\,d\pi.
 \label{eq:og-entropy}
\end{equation}
For ordinary reversal, Girsanov's formula gives the relative entropy of the
forward and reverse stationary hidden path laws on an interval of length
$T$ as $T\sigma$: their drift difference is $2v$ and their covariance is
$2D$. Equation~\eqref{eq:og-entropy} is dimensionless entropy per unit
physical time. A conversion into dissipated heat or power needs a valid
thermodynamic embedding and its temperature convention.

Let $f\in C^\infty(M;\mathbb R)$ be a deterministic noninvasive detector and
$Y_t=f(X_t)$. Define the observed local diffusivity
\begin{equation}
 V_t=\Gamma_f(X_t),\qquad
 \Gamma_f=\nabla f^{\mathsf T}D\nabla f,
 \qquad U_t=(Y_t,V_t).
 \label{eq:og-geometry}
\end{equation}
The scalar semimartingale quadratic variation satisfies
\begin{equation}
 [Y]_t-[Y]_s=2\int_s^t V_r\,dr.
 \label{eq:og-qv}
\end{equation}
Since $V$ has continuous paths, its value is half the time derivative of
this quadratic variation, including at a fixed time such as zero. Thus
$V_t$ is determined by an ideal continuous record of $Y$ and its calibrated
clock; it need not be a function of $Y_t$ alone.

This observation is a pathwise limiting statement. The square of one
increment, a regression of squared increments given $Y_t$, and a measurement
blurred by exposure averaging are not interchangeable with $V_t$. The
theorem below cannot be applied to them without an additional access and
error argument.

\section{The conditional local-defect theorem}

All conditional quantities in this chapter are defined up to the stationary
law's null sets. For $h>0$ put
\begin{align}
 g_h(U_0)
 &=\frac{\mathbb E[Y_h+Y_{-h}-2Y_0\mid U_0]}{2h},\\
 K_h(U_0)
 &=\frac{\mathbb E[(Y_h-Y_0)(Y_{-h}-Y_0)\mid U_0]}{h^2}.
 \label{eq:og-conditional-finite-h}
\end{align}
The symmetric definition of $g_h$ is essential. In general,
$\mathbb E[Af\mid Y_0,V_0]$ need not vanish, so a forward conditional
drift cannot silently replace it.

\begin{theorem}[Observed-geometry dissipation bound]
\label{thm:og-local}
Under the preceding assumptions, $g_h$ and $K_h$ converge in $L^\infty$
as $h\downarrow0$, with limits
\begin{align}
 g(U_0)&=\mathbb E[Sf(X_0)\mid U_0],\\
 K(U_0)&=\mathbb E[(Sf(X_0))^2-(Af(X_0))^2\mid U_0].
 \label{eq:og-conditional-limits}
\end{align}
Set $d(U)=g(U)^2-K(U)$, and define its ratio by zero on $V=0$. Then
\begin{equation}
 \boxed{\displaystyle
 \sigma\ \geq\ \mathcal I(Y)
 :=\mathbb E\!\left[
       \mathbf 1_{\{V_0>0\}}\frac{[d(U_0)]_+}{V_0}
                    \right].}
 \label{eq:og-local-bound}
\end{equation}
Equivalently, the truncated bounds
\[
 \mathcal I_\varepsilon(Y)
 =\mathbb E\!\left[
   \mathbf 1_{\{V_0\geq\varepsilon\}}
   \frac{[d(U_0)]_+}{V_0}\right]\leq\sigma,
 \qquad \varepsilon>0,
\]
increase to $\mathcal I(Y)$ as $\varepsilon\downarrow0$.
The number $\mathcal I(Y)$ depends only on the full stationary continuous
path law of $Y$, with its time scale, and not on the particular compatible
hidden diffusion used to calculate it.
\end{theorem}

\begin{proof}
Let $P_h$ and $P_h^\dagger$ be the forward and reverse semigroups. The past
and future are conditionally independent given $X_0$. Consequently,
\begin{align*}
 \mathbb E[Y_h-Y_0\mid X_0]&=(P_hf-f)(X_0),\\
 \mathbb E[Y_{-h}-Y_0\mid X_0]&=(P_h^\dagger f-f)(X_0),\\
 \mathbb E[(Y_h-Y_0)(Y_{-h}-Y_0)\mid X_0]
 &=(P_hf-f)(X_0)(P_h^\dagger f-f)(X_0).
\end{align*}
For a smooth function on the compact manifold, the semigroup expansion is
uniform in the hidden state:
\[
 \frac{P_hf-f}{h}=Lf+O(h),\qquad
 \frac{P_h^\dagger f-f}{h}=L^\dagger f+O(h).
\]
For example, the first remainder has supremum norm at most
$h\|L^2f\|_\infty/2$, by the integrated semigroup identity and its
supremum-norm contraction. Conditional expectation is an
$L^\infty$ contraction. Averaging the two expansions and multiplying them
therefore gives~\eqref{eq:og-conditional-limits} in $L^\infty$.
No Markov assumption on $Y$ has been used.

Abbreviate $a=Sf(X_0)$ and $c=Af(X_0)$. The conditional defect has the exact
form
\begin{equation}
 d(U_0)=\mathbb E[c^2\mid U_0]
             -\operatorname{Var}(a\mid U_0).
 \label{eq:og-defect-fiber}
\end{equation}
In particular,
\[
 [d(U_0)]_+\leq\mathbb E[c^2\mid U_0].
\]
Pointwise Cauchy--Schwarz in the diffusion metric gives
\[
 c^2=(v\mathbin{\cdot}\nabla f)^2
       \leq (\nabla f^{\mathsf T}D\nabla f)
             (v^{\mathsf T}D^{-1}v)=V_0 e(X_0).
\]
Since $V_0$ is included in $U_0$,
\[
 [d(U_0)]_+\leq V_0\mathbb E[e(X_0)\mid U_0].
\]
Divide on $V_0\geq\varepsilon$ and integrate to obtain the truncated
bound. Monotone convergence gives~\eqref{eq:og-local-bound}. On $V_0=0$,
uniform positive definiteness of $D$ implies $\nabla f=0$ and hence $c=0$.
Equation~\eqref{eq:og-defect-fiber} then gives $[d]_+=0$ on this event;
the declared zero value of the ratio is consistent.

Finally, continuous observation determines quadratic variation as a limiting
path functional, and~\eqref{eq:og-qv} determines its continuous derivative.
One precise construction uses the sums of squared increments along dyadic
partitions on each bounded interval. For the present bounded smooth It\^o
coefficients, these converge uniformly almost surely to quadratic variation:
the usual martingale estimate gives a squared supremum error of order the
mesh, and the dyadic meshes are summable. Defining the functional to be zero
on the null set of nonconvergent paths makes it Borel measurable. Taking its
right derivative at zero recovers $V_0$. Thus the continuous path law fixes
the joint law of $U_0,Y_{-h},Y_0,Y_h$ and the integrated conditional quantities
above for every $h$. Their unique $L^\infty$ limits, followed by the
nonnegative truncated limits, have the same value in every compatible model.
Versions on null sets cannot change that value.
\end{proof}

Although the quadratic-variation construction uses a path neighborhood,
under every admissible model its limit equals the instantaneous function
$\Gamma_f(X_0)$. Conditioning on this ideal limit therefore does not spoil
the conditional independence step given $X_0$. Replacing that limit with a
finite-window estimate built from the same increments is a different
statistical procedure and can introduce dependence and selection bias.

\section{The exact sources of slack}

The theorem provides more than an inequality: it separates two losses of
information. On $V_0>0$ write
\[
 m(U_0)=\mathbb E[c^2\mid U_0],\qquad
 w(U_0)=\operatorname{Var}(a\mid U_0).
\]
Since $m-[m-w]_+=\min(m,w)$, the following decomposition is exact:
\begin{align}
 \sigma-\mathcal I(Y)
 &=\mathbb E[e\,\mathbf 1_{\{V_0=0\}}]
   +\mathbb E\!\left[
      \left(e-\frac{c^2}{V_0}\right)
           \mathbf 1_{\{V_0>0\}}\right] \notag\\
 &\quad+\mathbb E\!\left[
      \frac{\min(m,w)}{V_0}\mathbf 1_{\{V_0>0\}}\right].
 \label{eq:og-slack}
\end{align}
Every term is nonnegative and finite. The second term is dissipation in
directions not detected by $f$ in the local diffusion metric. The third
records the loss caused by variation of the reversible drift inside a
retained observation fiber. In particular, equality holds if and only if,
up to stationary null sets:
\begin{enumerate}
 \item $e=0$ wherever $V_0=0$;
 \item on $V_0>0$, $v$ is a scalar multiple of $D\nabla f$;
 \item on $V_0>0$, either $\mathbb E[c^2\mid U_0]=0$ or
       $\operatorname{Var}(a\mid U_0)=0$.
\end{enumerate}
The conditions concern a particular compatible model. Their satisfaction
by one model, together with a path-law invariant lower bound, establishes
attainment of the minimum over the compatible class.

\section{Coordinate invariance and increasing physical access}

\begin{proposition}[Invertible scalar recalibration]
\label{prop:og-recalibration}
Let $\psi$ be a smooth one-to-one scalar transformation on an open interval
containing the range of $f$, with a smooth inverse and nonzero derivative.
Then the observed process $Z_t=\psi(Y_t)$ has
$\mathcal I(Z)=\mathcal I(Y)$.
\end{proposition}

\begin{proof}
Write $p=\psi'(Y_0)$ and $r=\psi''(Y_0)$. Diffusion calculus gives
\[
 \Gamma_{\psi\circ f}=p^2V_0,\qquad
 S(\psi\circ f)=p a+rV_0,\qquad
 A(\psi\circ f)=p c.
\]
The conditioning coordinates $(\psi(Y_0),p^2V_0)$ determine, and are
determined by, $(Y_0,V_0)$. Thus
\[
 g_Z=pg+rV_0,\qquad
 K_Z=p^2K+2prV_0g+r^2V_0^2,
 \qquad g_Z^2-K_Z=p^2(g^2-K).
\]
The ratio in~\eqref{eq:og-local-bound} is unchanged pointwise, including
its zero convention. The full bounds are invariant; truncations using the
same numerical threshold on the two different diffusivities need not be.
\end{proof}

For a complementary access statement, let $\mathcal G\subseteq\mathcal H
\subseteq\sigma(X_0)$ be instantaneous conditioning sigma-fields that both
contain $V_0$. Define the same conditional defect using each field; assume
the required conditioning variables are justified by the retained access.
The tower property and total conditional variance give
\[
 d_{\mathcal G}
 =\mathbb E[d_{\mathcal H}\mid\mathcal G]
  -\operatorname{Var}\!\left(
       \mathbb E[a\mid\mathcal H]\mid\mathcal G\right).
\]
Convexity of the positive part and the measurability of $V_0$ imply
\begin{equation}
 \mathcal I_{\mathcal G}\leq\mathcal I_{\mathcal H}\leq\sigma.
 \label{eq:og-refinement}
\end{equation}
Retaining more valid instantaneous information therefore cannot weaken this
population certificate. This statement does not assume that the refined
conditioning variables can be learned accurately from a finite dataset.

An exact variational form makes the connection with experimental witnesses
explicit:
\begin{equation}
 \mathcal I(Y)=
 \sup_{q,w}\mathbb E\!\left[
 w(U_0)\{2q(U_0)g(U_0)-q(U_0)^2-K(U_0)\}\right],
 \label{eq:og-dual}
\end{equation}
where $q,w$ are bounded measurable functions, $w\geq0$, and $w(U_0)V_0\leq1$.
For fixed $w$, the identity
$2qg-q^2-K=d-(q-g)^2$ shows that $q=g$ is optimal. Taking
$w=V_0^{-1}\mathbf 1_{\{V_0\geq\varepsilon,\,d>0\}}$ and letting
$\varepsilon\downarrow0$ proves equality. Each fixed admissible pair is
already a lower bound; the supremum is a population characterization, not
permission to optimize functions on the inference records without accounting
for selection.

\section{An attained minimum across hidden architectures}

Let $\mathbb P_{\gamma,\omega}^{\cos}$ denote the complete stationary
continuous path law of
\begin{equation}
 d\phi_t=\omega\,dt+\sqrt{2\gamma}\,dW_t\pmod{2\pi},\qquad
 Y_t=\cos\phi_t,
 \quad \phi_0\ \hbox{uniform},\quad \gamma>0.
 \label{eq:og-circle}
\end{equation}
Define $\mathfrak C(\mathbb P_{\gamma,\omega}^{\cos})$ to contain every
hidden model satisfying the smooth compact ordinary-diffusion assumptions of
this chapter, in any finite dimension, whose detector has precisely this
continuous path law with the same clock.

\begin{theorem}[Global compatible-model minimum for the cosine-circle law]
\label{thm:og-circle-minimum}
Within the declared class,
\begin{equation}
 \boxed{\displaystyle
 \min_{(X,f)\in\mathfrak C(\mathbb P_{\gamma,\omega}^{\cos})}
       \sigma(X)=\frac{\omega^2}{\gamma}.}
 \label{eq:og-circle-min}
\end{equation}
The circle model itself attains the minimum. Hidden dimension, topology,
and additional undetected coordinates are not fixed in this minimization.
\end{theorem}

\begin{proof}
For the circle, $Sf=-\gamma\cos\phi$, $Af=-\omega\sin\phi$, and
$\Gamma_f=\gamma\sin^2\phi$. Hence, with $y=Y_0$,
\[
 V_0=\gamma(1-y^2),\qquad
 g=-\gamma y,\qquad
 K=\gamma^2y^2-\omega^2(1-y^2).
\]
The stationary law of $Y_0$ is the arcsine law; in particular its endpoints
have probability zero. The ratio in~\eqref{eq:og-local-bound} is therefore
$\omega^2/\gamma$ almost surely. Path-law invariance implies this same
lower bound for every member of the compatible class. For the circle's
constant current velocity $v=\omega$ and diffusivity $D=\gamma$,
equation~\eqref{eq:og-entropy} gives $\sigma=\omega^2/\gamma$.
This feasible model attains the lower bound, proving the minimum.
\end{proof}

Full path-law agreement is sufficient but stronger than necessary for this
particular minimum. Let the compatibility requirement instead fix the
stationary joint law of $(Y_0,V_0)$ and the two local conditional functions
$g$ and $K$ appearing in~\eqref{eq:og-conditional-limits}. Every compatible
model still has the same $\mathcal I(Y)$, and the reference circle is still
feasible, so the minimum over this larger class remains
$\omega^2/\gamma$. These local population functions are substantially more
information than a finite list of lagged moments; the statement does not
remove the cost of resolving or estimating them.

This minimum is determined by a reversible observed law. Reversing the
stationary circle process changes $\omega$ to $-\omega$; reflecting the
phase $\phi\mapsto-\phi$ makes the same change and leaves $\cos\phi$
unchanged. Thus the observed forward and backward path laws agree, and their
relative-entropy rate is zero even when~\eqref{eq:og-circle-min} is positive.
The result does not contradict data processing: it uses additional
restrictions on admissible hidden physical models and the diffusion geometry
encoded in the observed law. It does not assert that $Y$ is Markovian.

The order of operations matters even for this exact example. Near
$|Y|=1$, the separate weighted terms $g^2/V$ and $K/V$ need not be
integrable; for the cosine-circle law their expectations are both positive
infinity. Their pointwise difference is nevertheless finite and equals
$\omega^2/\gamma$. Equation~\eqref{eq:og-local-bound} first forms the
conditional local difference, then its positive part, and then the
truncated integral. It is not a subtraction of divergent global integrals.

\subsection{A family with nonuniform stationary density and diffusivity}

The exact minimum is not confined to a constant-speed circle. Let
$\rho(\phi)$ and $D(\phi)$ be positive smooth even $2\pi$-periodic
functions, with $\int_0^{2\pi}\rho(\phi)\,d\phi=1$, and let $J$ be a
constant stationary probability flux. Define the circle diffusion by
\begin{equation}
 S=\rho^{-1}\partial_\phi(\rho D\partial_\phi),\qquad
 A=\frac{J}{\rho}\partial_\phi,
 \qquad Y=\cos\phi.
 \label{eq:og-symmetric-circle-family}
\end{equation}
For the complete observed law generated by this model, the minimum over all
compatible hidden diffusions in the chapter's class is
\begin{equation}
 \min\sigma
 =J^2\int_0^{2\pi}\frac{d\phi}{\rho(\phi)D(\phi)}.
 \label{eq:og-family-minimum}
\end{equation}
Indeed, $Sf$ is even, as are $\Gamma_f=D\sin^2\phi$ and
$e=J^2/(\rho^2D)$. Every even function on the circle is a measurable
function of $\cos\phi$. Thus $\operatorname{Var}(Sf\mid Y,V)=0$.
Away from the two detector critical points, which have zero stationary
probability, the one-dimensional Cauchy--Schwarz step is equality.
Equation~\eqref{eq:og-slack} proves $\mathcal I(Y)=\sigma$ for this model;
path-law invariance and the model's feasibility prove
\eqref{eq:og-family-minimum} across the compatible class. Reflection of the
phase preserves the invariant density, commutes with $S$, and changes $A$
into $-A$. It therefore maps the forward hidden generator to its stationary
reverse while leaving $Y$ fixed. Every scalar observed law in this family
is reversible, including the positive-minimum cases $J\ne0$.

More generally, a one-dimensional admissible hidden diffusion attains the
bound whenever $Sf$ is measurable with respect to $(f,\Gamma_f)$ and
$\nabla f\ne0$ almost surely. These are sufficient conditions to verify
in a model, rather than a claim that an arbitrary detector resolves its
hidden drift. Nonuniform diffusion or stationary density alone does not
prevent exact inference at the level of a full path law.

Equality of the full scalar covariance is a weaker specification than the
path law used here. For example, the earlier hub construction matching an
oscillatory covariance is a jump process and does not reproduce the
cosine-circle law's continuous paths and local diffusivity. Conversely, this
chapter gives no theorem that a fitted oscillatory covariance forces the
minimum~\eqref{eq:og-circle-min}.

For reference, the circle's finite-$h$ conditional quantities are exactly
\begin{align}
 g_h(y)&=\frac{e^{-\gamma h}\cos(\omega h)-1}{h}\,y,\\
 K_h(y)&=\frac{
 [e^{-\gamma h}\cos(\omega h)-1]^2y^2
 -e^{-2\gamma h}\sin^2(\omega h)(1-y^2)}{h^2}.
 \label{eq:og-circle-finite}
\end{align}
The untruncated finite-$h$ population ratio integrates to
\[
 \frac{e^{-2\gamma h}\sin^2(\omega h)}{\gamma h^2}
 \leq\frac{\omega^2}{\gamma},
\]
with equality in the limit $h\downarrow0$. This one-model identity does not
license an uncorrected finite-$h$ bound for arbitrary compatible diffusions.
A finite observation experiment needs explicit approximation, detector,
selection, and sampling guarantees.

\section{A vector extension}

The same argument has a matrix form. Let $F=(f_1,\ldots,f_k)$ be a smooth
vector detector and let
\[
 G_{ij}(X_0)=\nabla f_i^{\mathsf T}D\nabla f_j,
 \qquad U=(F(X_0),G(X_0)).
\]
The matrix $G$ is half the time derivative of the observed quadratic
covariation matrix. Put $a=SF$, $c=AF$, $g=\mathbb E[a\mid U]$, and
\[
 K_{\rm sym}
 =\lim_{h\downarrow0}\frac{1}{2h^2}\mathbb E\!\left[
 (F_h-F_0)(F_{-h}-F_0)^{\mathsf T}
 +(F_{-h}-F_0)(F_h-F_0)^{\mathsf T}\mid U\right].
\]
Uniform componentwise semigroup expansion gives
$K_{\rm sym}=\mathbb E[aa^{\mathsf T}-cc^{\mathsf T}\mid U]$.
If $G^+$ is the Moore--Penrose inverse and $H_+$ denotes the positive
spectral part of a symmetric matrix, then
\begin{equation}
 \sigma\geq
 \mathbb E\operatorname{tr}\left[
  \left((G^+)^{1/2}(gg^{\mathsf T}-K_{\rm sym})
                    (G^+)^{1/2}\right)_+\right].
 \label{eq:og-vector}
\end{equation}
This includes singular $G$; it is an inequality, not a claim of general
coordinate invariance under transformations of a singular vector detector.

To prove it, let $J$ be the detector differential, so that
$G=JDJ^{\mathsf T}$ and $c=Jv$. Linear algebra gives
$c^{\mathsf T}G^+c\leq e$ and $c\in\operatorname{range}(G)$.
Conditionally,
\[
 gg^{\mathsf T}-K_{\rm sym}
 =\mathbb E[cc^{\mathsf T}\mid U]
                -\operatorname{Cov}(a\mid U)
 \preceq\mathbb E[cc^{\mathsf T}\mid U].
\]
After whitening by $(G^+)^{1/2}$, the right side is positive semidefinite.
For symmetric $H\preceq N$ with $N\succeq0$, the variational identity
$\operatorname{tr}H_+=\max_{0\preceq P\preceq I}\operatorname{tr}(PH)$
gives $\operatorname{tr}H_+\leq\operatorname{tr}N$. Conditional expectation
and the preceding pointwise metric inequality establish
\eqref{eq:og-vector}. Truncating to events where $\|G^+\|$ is bounded
justifies every division; monotone convergence extends to their union,
which includes every finite matrix. As with the scalar theorem, a physically
available extra detector has an acquisition cost that must be recorded.

\section{Boundary of the result}

The minimum in~\eqref{eq:og-circle-min} is global over a broad, explicit
diffusion class, rather than over all conceivable physical descriptions.
The proof requires stationarity, smooth continuous deterministic detection,
the ordinary reversal convention for the complete hidden state, calibrated
time, and enough path access to determine local quadratic variation. It
does not cover arbitrary jump processes, singular latent dynamics,
generalized odd-variable reversal, or hidden sensor dynamics. Conditioning
on an averaged local diffusivity does not establish the required pointwise
geometry. No finite record proves exact path-law equality, and no laboratory
validation is supplied by this mathematical argument. These restrictions
are part of the theorem and are needed when deciding which application or
extension is scientifically justified.
